\documentclass[journal]{IEEEtran}
% IEEE TED submission draft (English). Build: paper/ted/build.sh (Tectonic). Figures: ../fig (paper/make_*figs*.py), ../../docs/fig/paper.
\usepackage{fontspec}
\setmainfont{texgyretermes}[Extension=.otf, UprightFont=*-regular, BoldFont=*-bold, ItalicFont=*-italic, BoldItalicFont=*-bolditalic]
\usepackage{amsmath}
\usepackage{graphicx}
\usepackage{booktabs}
\usepackage{cite}
\usepackage{array}
\usepackage{url}
\graphicspath{{../fig/}{../../docs/fig/paper/}}
\renewcommand{\topfraction}{0.92}\renewcommand{\textfraction}{0.06}\renewcommand{\floatpagefraction}{0.8}
\renewcommand{\dbltopfraction}{0.92}\renewcommand{\dblfloatpagefraction}{0.85}
\setcounter{topnumber}{3}\setcounter{dbltopnumber}{2}\setcounter{totalnumber}{5}

\begin{document}
\title{Tunable Inter-Body Hole Coupling in a Divided-Body SOI Single-Transistor Latch: Series-STL Equivalence, Continuous Tuning, and a Body Hole-Balance Model}

\author{Yeseul~Seo%
\thanks{Y. Seo is with the School of Electrical Engineering, Korea Advanced Institute of Science and Technology (KAIST), Daejeon, Republic of Korea (e-mail: ashleysys04@gmail.com).}%
\thanks{Simulation decks, analysis code and key logs are available at https://github.com/ashleysys04-blip/EE761\_DividedBody.}}

\markboth{IEEE Transactions on Electron Devices}{Seo: Tunable Inter-Body Hole Coupling in a Divided-Body SOI STL}
\maketitle

\begin{abstract}
A divided-body silicon-on-insulator (SOI) single-transistor latch (STL), in which a resistively grounded n$^+$ island splits the floating body into two series STLs inside one device, is studied with two-dimensional drift-diffusion TCAD. With hole transport across the island suppressed, the device latches twice, but it then reproduces two separately simulated STLs wired in series within 25~mV. Without lifetime killing, holes cross the island and couple the bodies: the second latch drops from 4.78 to 1.97~V, and the left STL departs from the series-circuit characteristic by up to 1.29~V with negative differential resistance. The island lifetime or width (15--200~nm) tunes the coupling continuously: with the first latch fixed at 1.44--1.45~V, the second latch moves between 1.97 and 4.78~V and the second-stage hysteresis between 0 and 2.27~V. A two-body hole-balance model explains the coupling. The holes are back-injected by the right STL into the island, in proportion to its body hole density and up to the plasma limit $\mu_p/(\mu_n+\mu_p)$, and cross the island by diffusion and drift over a decay length of 21~nm, derived from the island lifetime, the hole diffusivity and the loss into the via under the island. The model predicts the second-latch turn-off of every coupled island within 0.08~V and the turn-on of 30--90~nm islands within 0.2~V. A third jump, a regenerative n-p-n-p-n hole loop, was predicted for islands up to 40~nm only and confirmed by TCAD.
\end{abstract}

\begin{IEEEkeywords}
Single-transistor latch (STL), silicon-on-insulator (SOI), floating body, impact ionization, parasitic bipolar transistor, multi-state latch, hole coupling, compact model, carrier lifetime, TCAD.
\end{IEEEkeywords}

\section{Introduction}
\IEEEPARstart{I}{n} a floating-body SOI MOSFET, holes generated by impact ionization (II) at the drain junction charge the body, forward-bias the body--source junction and turn on the parasitic lateral n-p-n bipolar transistor, whose electron injection increases II again. When this loop gain exceeds unity the drain current jumps by several decades and returns at a lower voltage: the single-transistor latch (STL)~\cite{chen_stl}. STLs have been used as bistable resistors (biristors)~\cite{han_biristor}, capacitorless 1T-DRAM~\cite{kim_1tdram,moon_fin}, leaky integrate-and-fire neurons~\cite{han_lif,han_mimic,han_sciadv} and oscillators~\cite{kim_osc}, and low-voltage operation has been pursued~\cite{lee_1v}. These applications use two states. Ternary logic and multi-level elements need a third state; CMOS-compatible approaches include tunneling-based T-CMOS~\cite{jeong_tcmos}, a threshold switch in series with a MOSFET~\cite{heo_nbo2,han_bts,han_ptmos}, feedback FETs~\cite{son_fbfet} and gate--body diode MOSFETs~\cite{kim_spti}. A biristor threshold switch in series with a MOSFET~\cite{han_bts,han_ptmos} keeps leakage low but needs two devices. Recently a bulk MOSFET with two latches, the triristor, was demonstrated~\cite{noh_triristor}; its two latches share one p-well and are separated only by two current paths, so they can merge with bias.

A natural SOI alternative is to split the floating body physically. An n$^+$ island across the film divides the body into a left (source--island) and a right (island--drain) STL, and grounding the island through $R_{tap}$ lets the right STL latch first and then raise the island potential until the left STL latches. This paper first shows how such a device reaches two separate latches and why the simpler structures latch only once. It then asks the question a reviewer of any such device should ask: is a fully decoupled divided-body device anything more than two STLs in series? We show that it is not---it reproduces a series STL pair---and that what a single device adds is the \textit{hole coupling} between the two bodies across the island. That coupling makes the device depart from any series circuit, and its strength can be set continuously by the island lifetime or width.

The contributions are: (1) a step-by-step account of why body-doping splits, thickness splits, split gates, floating and tapped islands give one latch, and a hole-balance explanation of the two latches of the decoupled device (Sections~III--IV); (2) a direct comparison with two separately simulated STLs in series, and three signatures of the coupled device that a series circuit cannot produce (Section~V); (3) continuous tuning of the second latch and of the second-stage hysteresis by two knobs that collapse onto $L_p/W_{IS}$, with a gate-length control and the coupling limit (Section~VI); (4) a two-body hole-balance model in the equivalent-circuit form of~\cite{park_stl}, in which the coupling is derived (back-injection by the right STL and drift--diffusion across the island with a 21~nm decay length), that predicts the second-latch turn-off and turn-on and the onset of a third, regenerative jump (Section~VII).

\begin{figure*}[!t]
\centering
\includegraphics[width=\textwidth]{figW_concept.pdf}
\caption[Proposed divided-body SOI STL]{Divided-body SOI STL. (a) Left and right STLs share one gate; an n$^+$ island with a BOX via is grounded through $R_{tap}$; its width $W_{IS}$ and lifetime $\tau_{0,isl}$ set the hole coupling between the bodies. (b) Cross section of the reference device of Sections~IV--VII (dimensions in nm, body doping in cm$^{-3}$). (c) $|I_D|$--$V_D$ (solid: up, dashed: down) with the island killed ($\tau_{0,isl}=10^{-12}$~s) and coupled ($10^{-8}$~s, $W_{IS}$ = 30 and 120~nm); the first latch is common, the second moves from 4.78~V to 1.97~V with coupling. (d)--(g) Body holes in states I, II and III for the killed island (series pair: the left STL is turned on by the island potential $V_I$) and the coupled island (holes back-injected into the island diffuse on and charge the left body).
Data (draft only): (c) \path{1007_ddsplit/DDS_TE8_XS60.log_tr.log}, \path{1010_sink/NOKILL_TE8_XS60.log_tr.log}, \path{1030_islw/ISLW120.log_tr.log}; geometry \path{1007_ddsplit/DDS_TE8_XS60.in}; plot \path{paper/make_concept_figs.py}.}
\label{fig:concept}
\end{figure*}

\section{Device and Simulation Method}
The reference device (Fig.~\ref{fig:concept}) has $L_g=300$~nm, $T_{si}=50$~nm, $t_{ox}=20$~nm, $T_{box}=200$~nm, n$^+$ S/D of $10^{20}$~cm$^{-3}$, a 30~nm n$^+$ island ($10^{20}$~cm$^{-3}$) at 60\% of $L_g$ (left and right bodies 165 and 105~nm), left and right p-bodies of $6\times10^{17}$ and $7\times10^{17}$~cm$^{-3}$, $\tau_{0,Si}=10^{-8}$~s, $V_G=-0.2$~V, $R_{tap}=10$~M$\Omega$ and a width of 0.5~$\mu$m; the first devices of Sections~III and IV had the island at mid-channel (135/135~nm). The island is connected through an n$^+$ via in the BOX to a bottom tap electrode. Fig.~\ref{fig:circuit} gives the equivalent circuits: each STL is a MOSFET with a parasitic n-p-n transistor on its floating body, charged by an II current source; the island node is the source of the right STL and the drain of the left STL. With the island killed, the two STLs share only the island node [Fig.~\ref{fig:circuit}(a)]; with a coupled island, holes back-injected by the right STL cross the island and act as a hole source between the bodies [Fig.~\ref{fig:circuit}(b)].

Silvaco ATLAS~\cite{atlas} is used with \texttt{bgn}, \texttt{consrh}, \texttt{conmob}, \texttt{fldmob}, Fermi statistics, trap-assisted and band-to-band tunneling and the Selberherr II model~\cite{selberherr,maes}. With \texttt{consrh} the effective SRH lifetime is $\tau_{eff}=\tau_0/(1+N/N_{SRH})$, $N_{SRH}=5\times10^{16}$~cm$^{-3}$; lifetimes quoted below are the base value $\tau_0$, so $\tau_{eff}=\tau_0/2001$ in the island and $\approx\tau_0/14$ in the bodies. Because latching produces an S-shaped I--V on which a static voltage sweep stops at the fold, the drain is ramped in time, 0$\rightarrow$6$\rightarrow$0~V at 0.7~V/ms; the transient crosses the unstable branch and returns the jumps and hysteresis directly. A 100-fold change of the ramp rate moves the latches by $\le0.2$~V. The final runs use MixedMode (device plus $R_{tap}$ and sources), which agreed with device-only runs to four digits and was 20 times faster. A latch is defined as a change of the path current by more than one decade at a fixed applied voltage (within 2~mV) that reaches a real current level; numerical artifacts are excluded by checking the time step, the KCL residual, internal potentials and the gate displacement current.

\begin{figure}[!t]
\centering
\includegraphics[width=\columnwidth]{figX_circuits.pdf}
\caption[Equivalent circuits]{Equivalent circuits. Each STL is a MOSFET $M$ with a parasitic n-p-n $Q$ on its floating body $B$, charged by impact ionization $I_{ii}$; the island node $V_I$ is the drain of the left STL and the source of the right STL and is grounded through $R_{tap}$. (a) Killed island: the holes that the right STL back-injects into the island recombine there, so the two STLs interact only through $V_I$ (a series pair, Section~V). (b) Coupled island: part of them reaches the left body, a hole source $\eta I_D$ between the bodies, (\ref{eq:eta})--(\ref{eq:L}). Plot: \path{paper/make_concept_figs.py}.}
\label{fig:circuit}
\end{figure}

\section{Why the Device Latched Only Once}
\textit{Body-doping split.} Small (4--12\%) and large (2.3--10 times) left/right doping differences in a single body always gave one latch; the hole quasi-Fermi level stayed continuous across the doping step, because a p/p$^+$ step does not block majority holes, so the two regions cannot float to different potentials. An oxide notch under the body with a thin Si bridge, a left/right $T_{si}$ split, a split gate and an n$^+$/n$^-$/p/n$^+$ stack also gave one latch. \textit{Floating island.} A 30~nm n$^+$ island separated the two body potentials before latching ($-0.02$ and 0.36~V), but both bodies and the island jumped together at 3.26~V: with a floating island the right-STL current must flow through the left STL (series constraint). \textit{Tapped island.} Grounding the island through $R_{tap}$ removed this constraint, but at $R_{tap}$ = 600, 60 and 10~k$\Omega$ the latch stayed at 2.990~V; zooming into the jump, the right body rose first and the left body followed within $\approx10$~ns. Holes injected from the right body into the 30~nm island crossed it and charged the left body: the right p-body, the island and the left p-body form a lateral p-n$^+$-p transistor. A 100~nm island still coupled the bodies, and a top ohmic contact on the island let holes pass underneath it. Only reducing the lifetime of the island and via to 1~ps left the left STL off while the right STL latched alone. Raising $R_{tap}$ to 10~M$\Omega$ then lifted the island potential enough to drive the left STL, and a standalone left STL showed that a body doping of $6\times10^{17}$~cm$^{-3}$ is needed for an abrupt latch ($3$ and $5\times10^{17}$ turned on gradually).


\section{Decoupled Device: Two Latches and Their Mechanism}
With the island killed (1~ps), the reference device latches at 2.251~V (right STL; $I_{tap}$ jumps from $10^{-12}$ to $7\times10^{-8}$~A, state I$\rightarrow$II) and at 3.300~V (left STL; II$\rightarrow$III), when the island potential reaches 1.67~V, close to the latch voltage of the standalone left STL in the same model (1.78~V). On the way down the left and right STLs turn off at $\approx2.24$ and $\approx1.20$~V, so both latches have hysteresis.

Twenty-eight solution snapshots were analyzed per STL: source--body electron barrier, II hole current $I_{ii}=qW\!\int G_{ii}dA$ and recombination in the body and in the n$^+$ region that acts as its source. For the first latch the right drain-junction field is $1.16\times10^6$~V/cm ($M-1\approx0.25$); between 2.24 and 2.26~V, $I_{ii}$ jumps by four decades and the right barrier collapses from 0.60 to 0.28~eV. After latching, $I_{ii}$ equals body plus island recombination (e.g., $1.96\times10^{-8}\approx1.06\times10^{-8}+1.00\times10^{-8}$~A at 3.0~V). For the second latch the field at the left side of the island rises from $3.6\times10^5$ to $7.1\times10^5$~V/cm as the island potential rises, and between 3.28 and 3.32~V the left barrier collapses from 0.61 to 0.10~eV. The $G_{ii}$ map shows generation only at the two drain-side junctions; holes from the right body recombine in the island, so the left STL is turned on only by its own drain field. The turn-off is set by the hole margin: at $\tau_{0,Si}=10^{-7}$~s only 29\% of the generated holes recombine at the last on-point, so the latch survives to low current, whereas at $\tau_{0,Si}=10^{-8}$~s 78\% are lost and the left STL falls by six decades within 20~mV at a holding current of $5.5\times10^{-7}$~A.

Eighty-six single-parameter splits show that the two latches are designed independently: the right-body doping moves the first latch (1.24--2.51~V), the left-body doping the second (2.28--4.01~V), $V_G$ moves both, $R_{tap}$ divides the jump heights between them, and a shorter Si lifetime makes the turn-off abrupt at the cost of higher latch voltages. With $\tau_{0,Si}=10^{-8}$~s and the island moved to 60\% of $L_g$ (left body 165~nm, right body 105~nm), the two hysteresis loops lie at 1.23--1.45~V and 2.56--4.78~V and are separated by 1.12~V. This device (island 30~nm, island lifetime 1~ps) is the decoupled reference for the rest of the paper.



\section{Decoupled Device Equals a Series STL Pair}
To test whether the decoupled device adds anything to a circuit, the island of the decoupled reference was replaced by a metal core through the Si film with a 10~nm n$^+$ shell (an implementable form of island recombination; first and second latch 1.499/4.616~V), and this device was then cut through the middle of the core into two independent ATLAS devices, each keeping its geometry, mesh and half of the core as its internal electrode. The two STLs were wired in series in MixedMode with the same $R_{tap}$ and gate (CUT). A realistic pair of discrete SOI MOSFETs with 500~nm n$^+$ S/D and top contacts (SD) was also simulated. CUT reproduces the one-device result within 3~mV at the second latch and 25~mV elsewhere (Table~\ref{tab:ctrl}, Fig.~\ref{fig:series}(a)); SD differs by 0.19~V only through its geometry. A divided-body device with no hole transport across the island is therefore exactly a series STL pair: the separate latches alone are not a new function, only a more compact layout.

Keeping the island lifetime at the Si value ($\tau_{0,isl}=10^{-8}$~s, ``coupled'') changes this. The second latch drops by 2.6~V to 1.969~V, and the device departs from the series circuit in three ways. First, in a series circuit each STL current is set by its own terminal voltages, so the left-STL operating point ($|I_S|$, $V_{tap}$) must lie on the I--V of a standalone left STL with the same body (Fig.~\ref{fig:series}(b), gray). The killed-island device follows that curve within 0.03~V at 1--3~mA (2.269 vs 2.272~V at 1~mA). The coupled device carries the same current at a much lower voltage (1.724~V at 1~mA; 1.519 vs 2.469~V at 3~mA). Second, above 1~mA its left-STL voltage decreases as the current increases, a negative differential resistance that a monotonic standalone on-branch cannot produce. Third, a jump appears at 3.38~V. A hole balance of each body from the snapshots (generation, recombination and $\int J_p\,dy$ at the island edges) shows that the fraction of the left-body hole supply that has crossed the island rises from 23\% at 3.0~V to 59\% at 3.5~V: the jump closes an n-p-n-p-n hole loop from the right body through the island into the left body. With the island killed this fraction is zero and the jump is absent. Its condition is derived in Section~\ref{sec:model}. In circuit terms (Fig.~\ref{fig:circuit}(b)), the coupled device needs a body-to-body hole current source in addition to the $R_{tap}$ node; its magnitude is the coupling strength.

\begin{table}[!t]
\caption{Series-Circuit Control (V; Turn-Off at $|I_S|<10^{-8}$~A)}
\label{tab:ctrl}
\centering\footnotesize
\begin{tabular}{@{}lccc@{}}
\toprule
 & 1st on & 2nd on & 2nd off \\ \midrule
one device, metal core & 1.499 & 4.616 & 2.081 \\
two STLs in series (CUT) & 1.476 & \textbf{4.613} & 2.061 \\
two STLs in series (SD) & 1.471 & 4.428 & 2.019 \\
one device, coupled island & 1.444 & \textbf{1.969} & 1.747 \\ \bottomrule
\end{tabular}
\end{table}

\begin{figure*}[!t]
\centering
\includegraphics[width=0.82\textwidth]{figR_series.pdf}
\caption{Series-circuit control. (a) $|I_D|$--$V_D$ (solid: up, dashed: down) of the metal-core device, two separately simulated STLs in series (CUT, SD) and the coupled device; CUT overlaps the metal-core device. (b) Left-STL operating point ($V_{tap}$ vs $|I_S|$, down sweep) against a standalone left STL ($V_D$ vs $|I_D|$, gray). A series circuit must lie on the gray curve.}
\label{fig:series}
\end{figure*}

\section{Coupling Strength as a Design Knob}
Holes entering the n$^+$ island diffuse as minority carriers over $L_p=\sqrt{D_p\tau_{eff}}$; with $\tau_{eff}=\tau_{0,isl}/2001$ and $D_p\approx2$~cm$^2$/s, $L_p\approx32~\mathrm{nm}\times\sqrt{\tau_{0,isl}/10^{-8}~\mathrm{s}}$ (Section~VII derives 21~nm, with the simulator hole diffusivity and the loss into the via under the island; this rescales the $L_p/W_{IS}$ axis uniformly). The coupling should therefore scale with $L_p/W_{IS}$, and it was varied in two ways from the reference: (i) the island lifetime $\tau_{0,isl}=10^{-12}$--$10^{-8}$~s at $W_{IS}=30$~nm (only the two lifetime lines of the deck differ), and (ii) the island width $W_{IS}=15$--200~nm at $\tau_{0,isl}=\tau_{0,Si}=10^{-8}$~s. The width sweep keeps both body lengths and widens only the island, so $L_g=270~\mathrm{nm}+W_{IS}$ grows with it; a killed island at $W_{IS}=120$~nm separates this gate-length effect.

The latch windows are summarized in Fig.~\ref{fig:window}; I--V curves are shown in Fig.~\ref{fig:window}(a), (b) and Fig.~\ref{fig:model}(d), (e). The first latch (right STL) stays at 1.444--1.449~V in every case; only the second latch moves. In the lifetime sweep it stays at 4.74--4.78~V (series circuit) up to $\tau_{0,isl}=10^{-10}$~s and falls to 2.99~V at $10^{-9}$~s and 1.97~V at $10^{-8}$~s. In the width sweep it falls from 4.31~V (200~nm) through 3.34, 2.58, 2.19 and 2.06 to 1.97~V (30~nm). The two sweeps fall on one curve versus $L_p/W_{IS}$ [Fig.~\ref{fig:window}(c)--(e)]; at equal $L_p/W_{IS}$ they differ by at most $\approx0.3$~V, which we attribute to the one-dimensional diffusion picture behind $L_p/W_{IS}$. Killing the 120~nm island gives 4.761/2.577~V, the same as the killed 30~nm island (4.779/2.565~V), although $L_g$ grows from 0.30 to 0.39~$\mu$m; the width dependence of the coupled devices is therefore entirely due to coupling. The island width, a mask dimension, is the practical knob, whereas the lifetime knob needs a local, extreme lifetime control.

\begin{figure*}[!t]
\centering
\includegraphics[width=\textwidth]{figV_windows.pdf}
\caption[Latch windows]{Latch windows. (a), (b) $|I_D|$--$V_D$ of the width sweep ($\tau_{0,isl}=10^{-8}$~s) and the lifetime sweep ($W_{IS}=30$~nm); solid: up, dashed: down; filled/open circles: second latch-up $V_{LU,2}$ / latch-down $V_{LD,2}$; bar: one second-stage window. The 1~ns run has no down sweep. (c), (d) $V_{LU}$ (filled) and $V_{LD}$ (open) of the first (gray) and second (red) latches and the third jump (green) extracted from every run, versus island width at $\tau_{0,isl}=10^{-8}$~s and versus island lifetime at $W_{IS}=30$~nm; shaded: hysteresis windows; bar at 15~nm: continuous turn-on; thick lines: model of Section~\ref{sec:model} (solid $V_{LU,2}$, dashed $V_{LD,2}$); black lines in (c): killed 120~nm island; top axis of (d): $L/W_{IS}$, $L=21~\mathrm{nm}\times\sqrt{\tau_{0,isl}/10^{-8}~\mathrm{s}}$. (e) Second-stage hysteresis versus $L/W_{IS}$ for both knobs. (f) Jump heights of the second (red) and third (green) latches; band: predicted onset of the third jump.
Data (draft only; paths in the repository): I--V logs \path{1030_islw/ISLW{015,045,060,090,120,200,120_KILL}.log_tr.log}, \path{1103_third/ISLW{035,040}.log_tr.log}, \path{1010_sink/NOKILL_TE8_XS60.log_tr.log} (30~nm, 10~ns), \path{1029_isltau/ISL_TAU{1em11,1em10,1em9,3em9}_XS60.log_tr.log}, \path{1007_ddsplit/DDS_TE8_XS60.log_tr.log} (1~ps); extracted windows and jumps of (c)--(f): \path{paper/fig/figV_windows.dat}; third jump: \path{1103_model/third_onset.dat}; model: \path{1103_physics/stl_hybrid.py}; plot: \path{paper/make_window_fig.py}.}
\label{fig:window}
\end{figure*}


\textit{Hysteresis.} The second turn-off stays at 2.00--2.04~V for all coupled widths from 45 to 200~nm, so the width moves only the turn-on and the second-stage hysteresis is tuned continuously from 2.27~V to zero [Fig.~\ref{fig:window}(c), (e)]. The low turn-off is itself a coupling effect (2.04~V coupled vs 2.58~V killed at 120~nm): at 3.0~V on the down sweep the coupled island holds $10^{16}$--$10^{17}$~cm$^{-3}$ holes and keeps feeding the left body, whereas the killed island stays below $10^{4}$~cm$^{-3}$ (Section~\ref{sec:model}).

\textit{Limit of strong coupling.} The height of the second jump falls from 4.1 decades (200~nm) to 0.8 decades (30~nm), and at 15~nm the left STL turns on without a jump between 1.86 and 2.0~V and returns along the same curve: the two bodies behave as one hole reservoir. We take 30--120~nm as the tuning range and 200~nm, where the island is wider than either body, as the limit toward zero coupling.

\textit{Deviation from the series circuit.} With the operating-point test of Fig.~\ref{fig:series}(b) applied to every run, the left-STL voltage at 3~mA lies $-1.29$ to $-0.11$~V from the standalone characteristic for $W_{IS}$ = 15 to 200~nm and returns to $-0.03$/$+0.07$~V when the 30/120~nm islands are killed; the lifetime sweep lies on the same curve. The departure from a series STL pair thus holds over the whole tuning range, and its size is set by the coupling strength.


\section{Two-Body Hole-Balance Model}
\label{sec:model}
Following the equivalent-circuit approach of~\cite{park_stl}, each floating body is a node whose hole charge is set by the balance of hole generation and hole loss. The divided-body device needs one more element: a hole current source that links the two bodies through the island.

\subsection{Model Equations}
With $I$ the left-STL electron current $|I_S|$ and $V$ the island potential, the left body in steady state obeys
\begin{equation}
\bigl(M_L(V)-1\bigr)I+\eta\,I_D+G_0=\frac{I}{\beta_L(I)},\qquad I_D=I+\frac{V}{R_{tap}},
\label{eq:balance}
\end{equation}
where the left side collects impact-ionization holes, holes arriving from the island and other generation ($G_0=10^{-15}$~A), and the right side is the hole loss by recombination and back-injection into the source. In states II and III the right STL is latched, $M_R-1=1/\beta_R(I_D)$, which gives its on-branch $V_R(I_D)$, and $V_D=V+V_R(I_D)$. For every $I$, (\ref{eq:balance}) has one solution $V$ because its left side rises with $V$; the resulting $V_D(I)$ is S-shaped. The up sweep jumps at its local maxima, where the loop gain of~\cite{park_stl} reaches unity, and the down sweep drops at its local minima.

The coupling is the product of the holes the right STL injects into the island and the fraction that crosses it:
\begin{equation}
\eta=\gamma_R(I_D)\,T(W_{IS},\tau_{0,isl},I_D).
\label{eq:eta}
\end{equation}
\textit{Injection.} The island is the emitter of the right STL. The junction law $p_{isl}(0)\,n_{isl}=p_R\,n_R$ makes the back-injected hole current per drain electron proportional to the hole density $p_R$ of the right body. Extracted from the snapshots, $\gamma_R=0.019\,p_R/N_{A,R}$ holds for all widths up to $\approx$1~mA [Fig.~\ref{fig:model}(b)]; $p_R$ rises with current as the body is charged by its own impact ionization and then enters high injection. Above $\approx$10~mA the bodies and the island form a plasma ($n\approx p\approx10^{19}$--$10^{20}$~cm$^{-3}$) in which holes drift with the electrons, and $\gamma_R$ saturates at $\gamma_\infty=\mu_p/(\mu_n+\mu_p)$. We write
\begin{equation}
\gamma_R=\frac{g_R\,q_b}{1+g_R\,q_b/\gamma_\infty},\qquad q_b=\tfrac12\bigl(1+\sqrt{1+4I_D/I_{KF,R}}\bigr),
\label{eq:gamma}
\end{equation}
with the Gummel--Poon high-injection factor $q_b$.

\textit{Transmission.} Holes injected at the right island edge are collected at the left edge of an n$^+$ base of width $W_{IS}$. In the island they diffuse and drift in the ohmic field $E=J_D/(qN_D\mu_n)$ of the drain current: $L^2p''-L_Ep'-p=0$ with $L_E=\mu_pE\tau_p$, whose solution gives $T$ in closed form; $T\rightarrow1/\cosh(W_{IS}/L)$ at low current. The decay length follows from the island parameters alone. The bulk value $L_p=\sqrt{D_p\tau_p}$, with $\tau_p=\tau_{0,isl}/(1+N_D/N_{SRH})=5\times10^{-12}$~s and the hole diffusivity of the simulator in the $10^{20}$~cm$^{-3}$ island ($D_p=1.35$~cm$^2$/s), is 26.0~nm. Holes moving along the 50~nm film are also lost into the n$^+$ via under the island, a second sink deeper than $L_p$, so that $1/\tau_{eff}=(1+L_p/t_{si})/\tau_p$ and
\begin{equation}
L=\frac{L_p}{\sqrt{1+L_p/t_{si}}}=21.0~\mathrm{nm}.
\label{eq:L}
\end{equation}
\textit{Single-STL quantities.} $M-1$ and $1/\beta$ are extracted from the hole balance of the left body, as from a Gummel plot. For all runs, $M-1$ is one function of $V$ and $1/\beta$ one function of $I$; $1/\beta$ falls from 0.8 at $10^{-11}$~A to 0.05 near 1~$\mu$A and rises to 0.35 at 10~mA. The left STL's own impact-ionization holes generated on the island side of its drain junction ($\eta_{self}$, the $\eta$ of the killed island) are added to the generation. The island lifetime also changes the right STL, whose emitter it is: a killed island is a strong sink for its back-injection and needs 0.45~V more right-STL voltage at the same current, whereas coupled islands of any width give the same on-branch. Two reference curves $V_R(I_D)$ are therefore used, from the killed and the coupled 30~nm devices. Table~\ref{tab:model} lists the parameters; apart from these references and the 30~nm injection fit, every run is a prediction.

\begin{table}[!t]
\caption{Parameters of the Two-Body Hole-Balance Model}
\label{tab:model}
\centering\scriptsize\setlength{\tabcolsep}{3pt}
\begin{tabular}{@{}p{0.27\columnwidth}p{0.42\columnwidth}p{0.25\columnwidth}@{}}
\toprule
Quantity & Value / expression & Source \\ \midrule
$\tau_p$ (island) & $\tau_{0,isl}/(1+N_D/N_{SRH})$ & consrh \\
$D_p$, $\mu_n$ (island) & 1.35~cm$^2$/s, 90~cm$^2$/Vs & simulator \\
$L$ & $L_p/\sqrt{1+L_p/t_{si}}$ = 21.0~nm & (\ref{eq:L}) \\
$\gamma_\infty$ & $\mu_p/(\mu_n+\mu_p)$ = 0.37 & plasma limit \\
$T$ & drift--diffusion, $E=J_D/(qN_D\mu_n)$ & derived \\
$g_R$, $I_{KF,R}$ & 0.020, 7.4~$\mu$A & 30~nm injection \\
$M_L-1$ & $0.86/[1+e^{-(V-2.70)/0.52}]$ & extracted, all runs \\
$1/\beta_L$ & $0.032\,i^{-0.28}+0.016\,i^{0.50}/(1+(I/0.48\,\mathrm{mA})^{0.50})$, $i=I/1~\mu$A & extracted, all runs \\
$\eta_{self}$, $V_R(I_D)$ & killed / coupled 30~nm devices & references \\ \bottomrule
\end{tabular}
\end{table}

\subsection{Comparison with TCAD}
With the left STL off, the derived $\gamma_RT$ follows the extracted coupling over five decades from 15 to 200~nm and predicts the $\tau_{0,isl}=3\times10^{-9}$~s point [Fig.~\ref{fig:model}(a)]: the exponential width dependence and the $\sqrt{\tau}$ scaling of the two knobs come from one diffusion length, (\ref{eq:L}), which agrees with the 21.3~nm extracted from the coupling. Above 1~mA the drift term makes the island transparent ($T\rightarrow1$ for 15~nm), as in TCAD within 5--20\% [Fig.~\ref{fig:model}(c)]. Fig.~\ref{fig:model}(d)--(f) compares the predicted I--V. The second-latch turn-off is reproduced within 0.08~V for every coupled island (1.91--2.05 vs 1.96--2.04~V) and for the killed islands (2.58 vs 2.56--2.58~V): the 0.5~V lower turn-off of a coupled island is the island hole current that keeps the left body charged. On the on-branch, (\ref{eq:balance}) reads $M-1=1/\beta-\eta I_D/I$, so
\begin{equation}
\frac{d\ln I}{dV}=\frac{d\ln(M-1)/dV}{d\ln(1/\beta-\eta I_D/I)/d\ln I}:
\label{eq:slope}
\end{equation}
the slanted down sweep is the field sensitivity of impact ionization divided by the current roll-off of $\beta$. The turn-on is predicted within 0.09--0.18~V for 30--90~nm and exactly for $\tau_{0,isl}=3\times10^{-9}$~s (2.28~V); the 120 and 200~nm islands deviate by $-0.37$ and $+0.28$~V. The model does not reproduce the turn-on of killed islands, which occurs at island potentials above the extracted range of $M-1$ (2.8~V), nor the continuous turn-on at 15~nm, where it keeps a 1.5-decade jump.

\subsection{Third Jump: A Regenerative Hole Loop}
On the on-branch, the island holes replace part of the left body's own generation. When $\beta_L\eta\rightarrow1$ they replace almost all of it; $V(I)$ then decreases with current, and where this negative slope is steeper than the series resistance of the right STL, $V_D(I)$ folds a third time. The condition is that of a thyristor ($\alpha_{npn}+\alpha_{pnp}=1$) applied to the n-p-n-p-n loop right body--island--left body. At 50~mA, $\beta_L\eta$ extracted from the snapshots is 0.94, 0.88, 0.76 and 0.5 for 15, 30, 60 and 120~nm. Fitting the extracted $\eta(I_D,W_{IS})$ at each current over 15--120~nm gives $\eta=\eta_0(I_D)\exp[-W_{IS}/L_{eff}(I_D)]$ with $L_{eff}$ growing from 19 to 170~nm and $\eta_0\beta_L\rightarrow1.0$ above 10~mA; with this form the balance predicts the third fold for $W_{IS}\le40$~nm and none from 42.5~nm. Two additional TCAD runs confirmed the prediction: the jump within 1~mV is 0.61, 0.41, 0.29 and 0.21~decade at 15, 30, 35 and 40~nm and zero at 45, 60 and 90~nm [Fig.~\ref{fig:model}(f), green]. Across the jump the island hole density rises from $2$--$3\times10^{17}$ to $3.5$--$4.9\times10^{19}$~cm$^{-3}$, the plasma regime of (\ref{eq:gamma}). The present model, with $M-1$ independent of current, does not yet fold; the jump is an isothermal result at $\approx$6~mA/$\mu$m, and we do not claim it as a fourth state.

\begin{figure*}[!t]
\centering
\includegraphics[width=\textwidth]{figU_model.pdf}
\caption{Two-body hole-balance model. (a) Island coupling with the left STL off versus island width (circles), island lifetime (squares) and killed islands ($\times$), with the derived $\gamma_RT$ (\ref{eq:eta})--(\ref{eq:L}) and its prediction for $\tau_{0,isl}=3\times10^{-9}$~s (cross). (b) Holes back-injected by the right STL into the island per drain electron (symbols: 15--120~nm) and (\ref{eq:gamma}). (c) Island hole transmission, TCAD (symbols) and drift--diffusion (lines). (d), (e) $|I_D|$--$V_D$, TCAD (thin; solid up, dotted down) and model (thick; solid up, dashed down). (f) Second-latch turn-on and turn-off versus width, TCAD (filled) and model (open), and the third jump in TCAD (green, right axis) with the predicted onset (band).}
\label{fig:model}
\end{figure*}

\section{Discussion}
\textit{Implementation.} The two ends of the coupling range have different physical forms. Zero coupling needs an island effective lifetime of $\le5\times10^{-14}$~s or a metal core through the film (Section~V; modeled as an ideal ohmic boundary); intermediate coupling needs only the island width at the Si lifetime. The body lifetime of $\tau_{0,Si}=10^{-8}$~s used for abrupt turn-off is within conventional lifetime control (Au/Pt diffusion, electron or proton irradiation)~\cite{kim_lifetime}, and $R_{tap}$ = 1--10~M$\Omega$ can be a high-resistance poly-Si resistor.

\textit{Use of the coupling knob.} With the first latch fixed, the coupling sets where the second latch fires and how wide its hysteresis is. A wide second-stage window suits a three-state memory or a ternary threshold element; a narrow one gives a second, nearly reversible switching step. The decoupled device behaves as a series pair, so circuit models of series STLs, including current-driven relaxation oscillation, apply directly; for the coupled device the island node needs the additional hole current source $\eta I_D$ of (\ref{eq:balance})--(\ref{eq:L}), and the island lifetime also changes the right STL through its emitter loss.

\textit{Limitations.} All results are 2D, isothermal drift-diffusion. Local-field II overestimates generation in short channels: in the standalone left STL the energy-balance (hydrodynamic) latch was 0.7--0.8~V higher than the DD one. Full-device energy-balance runs of the decoupled device stopped at 3.6--5.3~V because the electron temperature at the right drain junction exceeded $10^4$~K and the time step collapsed, despite coupled Newton, a 1~nm mesh, relaxed tolerances, MixedMode and graded drains. An energy-balance run of the 30~nm coupled device, whose second latch lies lower, shows weaker II (no first latch below 3~V) and had not converged beyond 3.06~V; the second latch is therefore a DD result. An attempt to calibrate the II coefficients to a measured double-latch bulk device~\cite{noh_triristor} matched its Id--Vg but not its first latch, and a 20\% change of the Selberherr $B$ coefficient moved a coefficient-sensitive latch by $\approx$20\%. Absolute voltages are thus uncertain by roughly $\pm$20\%, whereas the comparisons in this paper (coupled vs killed vs series circuit at the same model) do not depend on the coefficients. Self-heating at the mA/$\mu$m currents of state III is not included; lattice-heating runs of the coupled device did not start in the present MixedMode setup.

\section{Conclusion}
A divided-body SOI STL with a resistively grounded n$^+$ island latches twice once hole transport across the island is removed, and each latch is the point where II hole generation in one body overcomes recombination. Such a decoupled device is, however, a series STL pair: two separately simulated STLs in series reproduce it within 25~mV. What a single device adds is the hole coupling between the bodies. The coupled device departs from the series characteristic by up to 1.29~V at 3~mA and shows negative differential resistance, and its coupling is tuned continuously by the island lifetime or, more practically, the island width. With the first latch fixed at 1.44--1.45~V, the second latch spans 1.97--4.78~V and the second-stage hysteresis 0--2.27~V, without a gate-length effect; the strongest coupling (15~nm) removes the second latch. A two-body hole-balance model shows where the coupling comes from: holes back-injected by the right STL into the island, in proportion to its body hole density and up to the plasma limit, that cross the island by diffusion over 21~nm, a length set by the island lifetime, the hole diffusivity and the via under the island, and by drift at high current. The model predicts the second-latch turn-off of all coupled islands within 0.08~V, the turn-on of 30--90~nm islands within 0.2~V and the slanted down sweep, and its regenerative-loop condition predicted that a third jump exists only for islands up to 40~nm, as TCAD confirmed. Remaining work is the current dependence of the multiplication and the off-state body bias, which set the turn-on of weakly coupled islands and the third fold, energy-balance and self-heating checks, and calibration against measured SOI devices.

\section*{Acknowledgment}
This work started as a course project of EE762 at KAIST.

\begin{thebibliography}{99}
\bibitem{chen_stl} C.-D. Chen, M. Matloubian, R. Sundaresan, B.-Y. Mao, C. C. Wei, and G. P. Pollack, ``Single-transistor latch in SOI MOSFETs,'' \emph{IEEE Electron Device Lett.}, vol.~9, no.~12, pp.~636--638, Dec. 1988.
\bibitem{han_biristor} J.-W. Han and Y.-K. Choi, ``Biristor---Bistable resistor based on a silicon nanowire,'' \emph{IEEE Electron Device Lett.}, vol.~31, no.~8, pp.~797--799, Aug. 2010.
\bibitem{kim_1tdram} D.-O. Kim, D.-I. Moon, and Y.-K. Choi, ``Optimization of bias schemes for long-term endurable 1T-DRAM through the use of the biristor mode operation,'' \emph{IEEE Electron Device Lett.}, vol.~35, no.~2, pp.~220--222, Feb. 2014.
\bibitem{moon_fin} D.-I. Moon, S.-J. Choi, J.-W. Han, S. Kim, and Y.-K. Choi, ``Fin-width dependence of BJT-based 1T-DRAM implemented on FinFET,'' \emph{IEEE Electron Device Lett.}, vol.~31, no.~9, pp.~909--911, Sep. 2010.
\bibitem{han_lif} J.-W. Han and M. Meyyappan, ``Leaky integrate-and-fire biristor neuron,'' \emph{IEEE Electron Device Lett.}, vol.~39, no.~9, pp.~1457--1460, Sep. 2018.
\bibitem{han_mimic} J.-K. Han \emph{et al.}, ``Mimicry of excitatory and inhibitory artificial neuron with leaky integrate-and-fire function by a single MOSFET,'' \emph{IEEE Electron Device Lett.}, vol.~41, no.~2, pp.~208--211, Feb. 2020.
\bibitem{han_sciadv} J.-K. Han \emph{et al.}, ``Cointegration of single-transistor neurons and synapses by nanoscale CMOS fabrication for highly scalable neuromorphic hardware,'' \emph{Sci. Adv.}, vol.~7, no.~32, Art. no. eabg8836, Aug. 2021.
\bibitem{kim_osc} H.-Y. Kim, S.-I. Kim, J.-K. Han, J.-W. Jung, and Y.-K. Choi, ``A single MOSFET-based oscillator on a bulk-silicon wafer,'' \emph{IEEE Electron Device Lett.}, vol.~45, no.~1, pp.~8--11, Jan. 2024.
\bibitem{lee_1v} S.-W. Lee \emph{et al.}, ``Single transistor latch near 1~V with asymmetric biasing in a MOSFET,'' \emph{IEEE Trans. Electron Devices}, vol.~71, no.~11, pp.~6539--6543, Nov. 2024.
\bibitem{han_ptmos} J.-A. Han, J.-W. Lee, J.-K. Han, D.-W. Kim, D.-H. Lee, and Y.-K. Choi, ``P-type ternary logic MOSFET with tunable middle state using bidirectional threshold switching,'' \emph{IEEE Electron Device Lett.}, vol.~46, no.~1, pp.~16--19, Jan. 2025.
\bibitem{han_bts} J.-K. Han, J.-M. Yu, S.-Y. Nam, and Y.-K. Choi, ``CMOS ternary logic with a biristor threshold switch for low static power consumption,'' \emph{IEEE Electron Device Lett.}, vol.~43, no.~7, pp.~1005--1008, Jul. 2022.
\bibitem{jeong_tcmos} J. W. Jeong \emph{et al.}, ``Tunnelling-based ternary metal--oxide--semiconductor technology,'' \emph{Nature Electron.}, vol.~2, no.~7, pp.~307--312, Jul. 2019.
\bibitem{heo_nbo2} S. Heo \emph{et al.}, ``High-speed ternary CMOS inverter by monolithic integration of NbO$_2$ threshold switch with MOSFET,'' in \emph{IEDM Tech. Dig.}, Dec. 2021, pp.~32.2.1--32.2.4.
\bibitem{son_fbfet} J. Son, K. Cho, and S. Kim, ``New ternary inverter with memory function using silicon feedback field-effect transistors,'' \emph{Sci. Rep.}, vol.~12, Art. no. 12907, Jul. 2022.
\bibitem{kim_spti} H.-Y. Kim, H.-B. Noh, H.-Y. Maeng, and Y.-K. Choi, ``A selective positive ternary inverter using the mirror symmetry of a MOSFET with a gate--body diode,'' \emph{IEEE Trans. Electron Devices}, vol.~73, no.~8, pp.~5232--5238, Aug. 2026.
\bibitem{park_stl} J.-H. Park, H.-B. Noh, S.-W. Lee, S.-Y. Yun, and Y.-K. Choi, ``Analytical model for single transistor latch in MOSFETs,'' \emph{IEEE Electron Device Lett.}, early access, 2026, doi: 10.1109/LED.2026.3737574.
\bibitem{noh_triristor} H.-B. Noh \emph{et al.}, ``A triristor with tri-stable states using double latches,'' \emph{IEEE Electron Device Lett.}, vol.~46, no.~8, pp.~1269--1272, Aug. 2025.
\bibitem{kim_lifetime} S. Kim, S.-J. Choi, D.-I. Moon, and Y.-K. Choi, ``Carrier lifetime engineering for floating-body cell memory,'' \emph{IEEE Trans. Electron Devices}, vol.~59, no.~2, pp.~367--373, Feb. 2012.
\bibitem{maes} W. Maes, K. De Meyer, and R. Van Overstraeten, ``Impact ionization in silicon: A review and update,'' \emph{Solid-State Electron.}, vol.~33, no.~6, pp.~705--718, Jun. 1990.
\bibitem{selberherr} S. Selberherr, \emph{Analysis and Simulation of Semiconductor Devices}. Vienna, Austria: Springer-Verlag, 1984.
\bibitem{atlas} Silvaco Inc., \emph{Atlas User's Manual}, version 5.32.1.R, Santa Clara, CA, USA.
\end{thebibliography}

\end{document}
